\begin{ouverture}[calcul]{3}{Découvrir les nombres relatifs}{Partie 1 · Nombres et calculs}
{Zéro · Opposé · Valeur absolue · Ordre}
\objectif{Lire et placer des nombres sur une droite.}
\objectif{Comparer des décimaux relatifs.}
\objectif{Distinguer signe et distance à zéro.}
\prerequis{Repérage et nombres décimaux.}
\end{ouverture}
\pagestyle{otMain}\essentiel\fontsize{11}{15}\selectfont
\phantomsection\label{cours-3-1}
\essSection{1}{De part et d’autre de zéro}
\begin{definition}Les nombres relatifs permettent de représenter des valeurs inférieures à zéro. $-4$ peut être une température de quatre degrés sous zéro. Un nombre positif est supérieur ou égal à 0 ; un nombre strictement positif est supérieur à 0. De même, négatif signifie inférieur ou égal à 0, strictement négatif signifie inférieur à 0. Zéro est à la fois positif et négatif.\end{definition}
\begin{definition}Deux nombres sont opposés lorsque leur somme est nulle : $-3$ et $3$ sont opposés car $(-3)+3=0$. Sur une droite graduée, ils sont à la même distance de zéro, de part et d’autre. L’opposé de $a$ se note $-a$ ; ainsi $-(-3)=3$ : $-a$ n’est pas toujours négatif. La valeur absolue d’un nombre est sa distance à zéro : $-3$ et $3$ ont tous deux pour valeur absolue 3. La notation $\lvert a\rvert$ n’est pas exigible en Cinquième.\end{definition}
\label{fig-cours-3-1-cours-relatifs}\figureGeo{cours-relatifs}

\begin{attention}Le signe placé devant le nombre fait partie de son écriture. Sur une droite, l’unité de graduation doit rester constante.\end{attention}
\clearpage
\phantomsection\label{cours-3-2}
\essSection{2}{Comparer et ranger}
\begin{propriete}Sur une droite orientée vers la droite, les nombres augmentent de gauche à droite. Tout nombre strictement négatif est inférieur à 0 et à tout nombre positif. Entre deux nombres négatifs, celui qui est le plus loin de zéro est le plus petit.\end{propriete}
\begin{exemple}$-7<-2$, alors que $-7$ est plus loin de zéro que $-2$ : sa valeur absolue est plus grande. Pour les décimaux, $-2{,}8<-2{,}35$ : $-2{,}8$ est plus loin de zéro que $-2{,}35$.\end{exemple}
\begin{savoirfaire}[Construire un ordre]Ranger d’abord les nombres négatifs, puis zéro, puis les positifs. Entre deux nombres, on peut toujours trouver un décimal : entre $-1{,}4$ et $-1{,}3$, on peut choisir $-1{,}35$.\end{savoirfaire}
\begin{savoirfaire}[Interpréter le signe]Indiquer la référence : niveau de la mer, température de fusion de l’eau, origine d’un axe. Le signe seul n’indique pas une distance parcourue.\end{savoirfaire}

\clearpage\rubrique{Découverte}
\begin{blocLivre}{ex-3-1}\begin{exercice}[1]{Lire une graduation}
Lire les abscisses de A, B et C sur la droite.
\label{fig-ex-3-1}\figureGeo{rel-1}
\end{exercice}\end{blocLivre}
\begin{blocLivre}{ex-3-2}\begin{exercice}[1]{Placer deux nombres}
Placer $D(-2)$ et $E(3)$ sur la droite donnée.
\label{fig-ex-3-2}\figureGeo{rel-2}
\end{exercice}\end{blocLivre}
\begin{blocLivre}{ex-3-3}\begin{exercice}[1]{Trouver les opposés}
Donner les opposés de $-8$, $2{,}4$ et 0.
\end{exercice}\end{blocLivre}
\begin{blocLivre}{ex-3-4}\begin{exercice}[1]{Distances à zéro}
Sur une droite graduée, à quelle distance de zéro sont les points d’abscisses $-6{,}5$ et 4 ? En déduire la valeur absolue de chacun de ces nombres.
\end{exercice}\end{blocLivre}
\begin{blocLivre}{ex-3-5}\begin{exercice}[1]{Un premier ordre}
Comparer $-5$ et 2, puis $-9$ et $-4$.
\end{exercice}\end{blocLivre}

\rubrique{Entraînement}
\begin{blocLivre}{ex-3-6}\begin{exercice}[2]{Ranger des décimaux}
Ranger dans l’ordre croissant, puis dans l’ordre décroissant : $1{,}2$ ; $-1{,}02$ ; 0 ; $-1{,}2$ ; $0{,}8$.
\end{exercice}\end{blocLivre}
\begin{blocLivre}{ex-3-7}\begin{exercice}[2]{Encadrer un négatif}
Encadrer $-3{,}7$ entre deux entiers consécutifs.
\end{exercice}\end{blocLivre}
\begin{blocLivre}{ex-3-8}\begin{exercice}[2]{Affiner une graduation}
Lire les abscisses de F et G. Une unité est partagée en quatre.
\label{fig-ex-3-8}\figureGeo{rel-8}
\end{exercice}\end{blocLivre}
\begin{blocLivre}{ex-3-9}\begin{exercice}[2]{Zéro et les signes}
Parmi $-2$, 0 et 5, lister les nombres positifs, puis strictement positifs.
\end{exercice}\end{blocLivre}
\begin{blocLivre}{ex-3-10}\begin{exercice}[2]{Même distance}
Sur une droite graduée, quels nombres sont situés à 4 unités de zéro ? Et à 0 unité de zéro ?
\end{exercice}\end{blocLivre}

\rubrique{Problèmes}
\begin{blocLivre}{ex-3-11}\begin{exercice}[3]{Comparer des altitudes}
Quatre lieux ont pour altitudes $-25$ m, 12 m, $-8$ m et 0 m. Les ranger du plus bas au plus haut et identifier ceux sous le niveau de la mer.
\end{exercice}\end{blocLivre}
\begin{blocLivre}{ex-3-12}\begin{exercice}[3]{Trouver un intermédiaire}
Proposer deux décimaux distincts strictement compris entre $-2{,}4$ et $-2{,}3$.
\end{exercice}\end{blocLivre}
\begin{blocLivre}{ex-3-13}\begin{exercice}[3]{Des contraintes entières}
Lister les entiers $x$ vérifiant $-3<x\le2$.
\end{exercice}\end{blocLivre}
\begin{blocLivre}{ex-3-14}\begin{exercice}[3]{Une fausse règle}
« Le nombre le plus éloigné de zéro est toujours le plus grand. » Donner un contre-exemple.
\end{exercice}\end{blocLivre}
\begin{blocLivre}{ex-3-15}\begin{exercice}[3]{Classement et opposés}
On sait que $-6<a<-2$. Sans connaître $a$, encadrer son opposé $-a$. Le nombre $-a$ est-il négatif ? Expliquer avec une droite.
\label{fig-ex-3-15}\figureGeo{rel-15}
\end{exercice}\end{blocLivre}

\clearpage
